\section{Related Work}
\label{sec:related}

We position this SoK against three prior literature clusters: surveys
and taxonomies of zero-knowledge proof systems, prior work on
privacy-preserving group protocols, and prior SoKs in adjacent
privacy-preserving primitive areas. The first cluster overlaps
directly with our taxonomic ambition; the second and third overlap
with the problem class.

\subsection{Surveys and taxonomies of zero-knowledge proof systems}

The cryptographic literature contains substantial survey work on
zero-knowledge proof systems. Treatments index variously by
construction family --- pairing-based versus IPA-based versus FRI-based
--- by setup posture --- per-circuit, universal-updatable, transparent
--- or by underlying hardness assumption.

What distinguishes the present work from these surveys is the
indexing axis. Existing surveys treat the proof system as the
primary object of interest and the deployment context as a
consequence; we invert this and treat the deployment constraint
(anchor chain, host functions, fee economics) as the primary axis,
with the proof system following as a derived choice. This inversion
is what makes the cross-axis dependency analysis of
\S\ref{sec:dimensions-deps} legible and what justifies the
profile-conditional recommendations of
\S\ref{sec:recommendations}. A purely cryptographic taxonomy would
not produce \S\ref{sec:rec-jurisdiction}'s ``none of C1--C6''
non-recommendation, because the cryptographic-taxonomy lens does not
weight operational constraints.

\paragraph{TODO §9 verification pass.} Lock specific citation
keys to: (i) the Nitulescu zk-SNARK introduction, (ii) the
B\"unz et al.\ Bulletproofs / IPA lineage, (iii) any
ePrint surveys post-2020 that index zkSNARKs by deployment
constraint. The §9 research pass is the
citation-heavy half of this section's work.

\subsection{Privacy-preserving group protocols}

The closest prior work to the registry abstraction of
\S\ref{sec:problem} is Signal's Private Group
System~\cite{chase-2020-signal-pgs}, which closes group-membership
metadata against an honest-but-curious application-layer operator
using anonymous credentials. The Signal construction differs from
the registry abstraction in two respects. First, it places an
application-layer operator in the trust path; the registry
abstraction's contribution is precisely to remove that operator.
Second, it does not anchor state on a public ledger; this is a
property the operator-removal trade is in exchange for. The two
constructions are complementary rather than competing: an MLS-shaped
deployment using Signal's credential machinery for membership and
a registry of the kind surveyed here for state anchoring would
combine their guarantees.

Tornado Cash~\cite{pertsev-2019-tornado} and
Semaphore~\cite{semaphore-2024} are direct instantiations of the
registry abstraction with single-relation, fixed-policy
configurations. They sit at C2 in the taxonomy and have driven the
substantial ecosystem of BN254-precompile-aware tooling that
configurations C2 and C5 inherit. Both pre-date the host-function
heterogeneity that makes this paper's taxonomy necessary.

\paragraph{TODO §9 verification pass.} Cite specific MLS architecture
papers (Beurdouche et al.\ MLS, the IETF MLS RFC) and the Coconut
threshold-credential lineage (Sonnino et al.) as the two-cluster
basis for ``privacy-preserving group protocols outside our
abstraction.''

\subsection{SoKs in adjacent primitive areas}

The SoK genre has been applied to several adjacent privacy-preserving
primitive classes: privacy-preserving payments and mixers,
anonymous credentials, decentralized identity, and mixnets. These
prior SoKs have established the methodology of comparing
constructions against an explicit threat model with profile-conditional
recommendations, which the present work follows. The methodological
debt is direct.

The substantive distinction from each:
\begin{itemize}
  \item Privacy-preserving-payments SoKs treat a different primitive
    --- value transfer with confidential amounts and unlinkable
    senders --- and frame the problem around mixer constructions and
    confidential-transaction protocols. Tornado Cash sits in both
    that primitive class and ours; the per-payment relation differs
    structurally from the per-state-transition relation we treat as
    canonical.
  \item Anonymous-credentials SoKs index along a different axis:
    issuer / prover / verifier roles, predicate expressivity, and
    revocation handling. The registry abstraction does not have an
    issuer in the credential sense.
  \item Mixnet SoKs target a different threat model:
    network-observer-only adversaries with active routing
    capability. The registry abstraction's threat model includes the
    chain-validator as a party, which mixnet literature
    typically excludes by construction.
  \item Decentralized-identity SoKs assemble the credential, registry,
    and revocation primitives into composite systems. The present
    work treats one of those primitives (the registry) in isolation.
\end{itemize}

\paragraph{TODO §9 verification pass.} Lock keys to specific SoKs in
each adjacent area. At minimum: a privacy-preserving-payments SoK,
an anonymous-credentials SoK, a mixnet SoK, and a
decentralized-identity SoK. The structural claims above are defensible
independently but would be strengthened by direct citations.

\subsection{Position}

Our contribution is the deployment-constraint-indexed taxonomy of
\S\ref{sec:dimensions}, the cross-axis dependency analysis of
\S\ref{sec:dimensions-deps}, and the profile-conditional
recommendations of \S\ref{sec:recommendations}. Prior work has either
(a) surveyed the cryptography without the deployment constraints, or
(b) constructed individual systems that instantiate one point in our
design space. We are not aware of a prior SoK that takes the
metadata-hiding-group-registry primitive as the unit of comparison
and indexes by deployment constraint. The present work occupies that
position.
